%%---------------------------------------------------------------------------%%

%  Introduction
%%%%%%%%%%%%%%%%%%%%%%%%%%%%%%%%%%%%%%%%%
\section{Introduction}

\subsection{Introduction to CASL}
The Consortium for Advanced Simulation of Light Water Reactors (CASL) is the first
US Department of Energy (DOE) Energy Innovation Hub, established in July 2010 for the purpose of providing
advanced modeling and simulation (ModSim) solutions for commercial nuclear reactors.

CASL's vision is to predict, with confidence, the performance of nuclear reactors
through comprehensive, science-based modeling and simulation technology that is
deployed and applied broadly throughout the nuclear energy industry to enhance safety,
reliability, and economics.

CASL's mission is to provide coupled, high-fidelity, usable modeling and simulation
capabilities needed to address light-water reactor (LWR) operational and safety performance-defining phenomena.

CASL's foundational technology products include CASL solutions and CASL ModSim Technologies.
CASL's ModSim technology, the Virtual Environment for Reactor Applications (VERA), provides
higher fidelity results than those offered by the current industry approach by incorporating coupled physics
and science-based models, state-of-the-art numerical methods, modern computational science,
integrated uncertainty quantification (UQ) and validation against data from operating
pressurized water reactors (PWRs), single-effect experiments, and integral tests.

Through new insights afforded by its ModSim technology, CASL will address key nuclear
energy industry challenges to increasing power uprates, ensuring higher fuel burnup, and extending lifetimes
while providing higher confidence in enhanced nuclear safety and this cleaner energy source.

The CASL team is a consortium that consists of ten core partners and numerous contributing members.
The CASL organization is led by Oak Ridge National Laboratory (ORNL), and CASL's research and development is executed
by six technical teams or \textit{focus areas} (FAs) and one integrating technical area.
This groundbreaking partnership provides unparalleled
collective institutional knowledge, nuclear science and engineering talent, computational
science leadership, and a record of LWR design and regulatory accomplishments.

More information on CASL can be found at the website \href{http://www.casl.gov}{www.casl.gov}.

%%%%%%%%%%%%
\subsection{VERA}
VERA is a specific collection
of multi-physics computer codes used to model and simulate depletion of an LWR core over multiple cycles.
Examples of the separate physics modeled in the core simulator include cross section generation,
neutron transport, isotopic depletion, thermal hydraulics, and fuel performance.

The purpose of the core simulator is to simulate depletion of the reactor core and provide data and boundary conditions to model
CASL Challenge Problems such as crud-induced power shift (CIPS), crud-induced localized corrosion (CILC), departure from nucleate boiling (DNB), pellet-cladding interaction (PCI), and reactivity insertion accident (RIA) analyses.

One important feature of the core simulator is that a single common input file is used to drive all
of the different physics codes.\footnote{The only exception to this is for computational fluid dynamics codes,
which generally require a detailed CAD file to support mesh generation and perform meaningful analysis.}
One benefit of using a single common input is that users only need to understand and be proficient with one input
instead of having to understand multiple inputs for multiple physics codes.
Another benefit of using a single common input is that all codes work from a single
geometry description, and this reduces errors due to inconsistent geometries in different codes.

The most up-to-date version of this document resides in the VERA Git repository file
``VERAIO/verain/docs/verain\_UM.pdf''.
Please refer to this location for the latest version of the input manual.


%%%%%%%%%%%%
\subsection{Manual Organization}

This manual is organized into three main parts.

The first part, which includes chapters~\ref{chap:user} through \ref{chap:depletion}, is the ``User's Manual,''
which describes how a user would set up a typical input. This part of the manual gives
the most common input cards a user would need and describes how to use them.
It does not include a complete list of cards or show every available option.

The second part of the manual, Chapter~\ref{chap:cards}, is a ``Reference Manual'' and includes
a complete list of every available input card.

The third part of the manual, Chapter \ref{chap:example}, gives several example input decks.
Additional example input files can be found in the code installation directory.

In addition, a description of the VERARun script that is used to run VERA jobs is given in Chapter~\ref{chap:verarun}.

Note that the VERA input processor, VERAIn, is an open-source software project and can be
found on the CASL Github website \href{http://www.github.com/casl/verain}{github.com/casl/verain}.
The open-source input processor does not include any physics packages.


%%%%%%%%%%%%
\subsection{Training Requirements}

There is no required training for running VERA, but users should have a basic understanding
of LWR technology. Users who perform any engineering or safety-related work with VERA should
follow the procedures of their own organizations.

Optional user training is periodically available from the VERA Users Group. 
Please contact support (contact information given below) to inquire about training opportunities.

%%%%%%%%%%%%
\subsection{Purpose and Functional Requirements}

The purpose and functional requirements of the VERA common input processor (VERAIn) are:
\begin{enumerate}
\item Read an ASCII input provided by the user as described in this manual.
\item Perform basic error checking on the ASCII input. Additional error checking is performed by
other VERA components.
\item Perform basic geometry processing such as expanding input maps from octant to full geometry
where applicable.
\item Create an extensible markup language (XML) output file that can easily be read by other VERA components.
\end{enumerate}

The purpose of the VERA run script (VERARun) is to provide a single
interface to run the VERA codes, usually in parallel computing environments.
The specific functional requirements for VERARun are:
\begin{enumerate}
\item Run VERAIn to create an XML file that can be read by other VERA components.
\item Run any input preprocessors as necessary (such as XML2CTF or XML2Bison).
\item Submit jobs to a parallel computing cluster. 
\end{enumerate}


%%%%%%%%%%%%
\subsection{Code Capabilities and Limitations}

The current code capabilities of VERAIn and VERARun are specified by the functional requirements listed above.
Requirements not explicitly stated in this list are assumed to be limitations.

One general limitation in the input processor is that the input is limited to standard LWR designs.
For example, the input processor does not support reactors with hexagonal or plate fuel or with coolant
that is not water. 

Other VERA components may have limitations; the user should refer to the documentation
of the other VERA components for these limitations.

%%%%%%%%%%%%
\subsection{Testing}

Information regarding system testing can be found in the respective VERAIn and VERARun Software Test Plan
Requirements and Test Reports. These documents contain summaries of all system testing and associated requirements.
Any feature not covered in these reports is considered to be untested.

%%%%%%%%%%%%
\subsection{Computer System Vulnerabilities}

Running VERAIn or VERARun on any machine is not known to expose the system to any security
vulnerabilities at this time.
VERAIn and VERARun should not be run with administrative level access permissions.

%%%%%%%%%%%%
\subsection{Software Support}

For specific questions about the use of VERAIn or VERARun, the licensing of the code, or to report bugs,
users should send an email to {\tt casl-support@ornl.gov}.

Additional user information may also be found on the CASL website \href{http://www.casl.gov}{www.casl.gov}.








